% 2. แนวคิดและงานที่เกี่ยวข้อง -- ย่อจากรายงานบทที่ 2 (../report/chapter/02-literature.tex)
\section{แนวคิดและงานที่เกี่ยวข้อง}

\subsection{สุขภาวะทางจิตและแบบวัด DASS-21}

งานวิจัยนี้แยกสุขภาวะทางจิตเป็นสามด้านตามโมเดลสามองค์ประกอบ (tripartite model)
ของ Clark และ Watson \cite{clark1991tripartite}
ซึ่งอธิบายว่าภาวะซึมเศร้ากับความวิตกกังวลมีส่วนที่ทับซ้อนกันคือความทุกข์ทั่วไป
และมีส่วนที่แยกจากกัน ภาวะซึมเศร้ามาคู่กับการขาดอารมณ์ด้านบวก
ส่วนความวิตกกังวลมาคู่กับการตื่นตัวทางกาย
ถ้าใช้คะแนนรวมค่าเดียว จะตอบไม่ได้ว่าโซเชียลมีเดียสัมพันธ์กับด้านใดมากกว่ากัน

แบบวัด DASS-21 \cite{lovibond1995manual} ถามถึงความรู้สึกในช่วง 1 สัปดาห์ที่ผ่านมา
มี 21 ข้อ แบ่งเป็นมิติภาวะซึมเศร้า ความวิตกกังวล และความเครียด มิติละ 7 ข้อ
แต่ละข้อให้คะแนน 0 ถึง 3 แล้วนำผลรวมรายมิติคูณ 2 จึงได้คะแนนมิติละ 0 ถึง 42
คะแนนยิ่งสูงแปลว่ามีอาการมาก
แบบวัดนี้มีฉบับแปลภาษาไทยบนเว็บไซต์ทางการของแบบวัด \cite{webster2008thai}
และผ่านการตรวจสอบคุณภาพในนักศึกษาไทยแล้ว \cite{wittayapun2023validation}

\subsection{โซเชียลมีเดียกับสุขภาวะทางจิต}

กลไกที่มักใช้อธิบายความสัมพันธ์มีสามทาง ทางแรกคือการแย่งเวลา (displacement)
คือเวลาหน้าจอไปเบียดเวลานอน การออกกำลังกาย และการพบปะผู้คน
ทางที่สองคือการเปรียบเทียบทางสังคม (social comparison) \cite{festinger1954social}
ผู้ใช้มักเทียบตัวเองกับภาพชีวิตของผู้อื่นที่ผ่านการคัดมาแล้ว ซึ่งสัมพันธ์กับความรู้สึกด้อยค่า
ทางที่สามคือ \textit{doomscrolling} หรือการไถดูข่าวเชิงลบไปเรื่อย~ๆ จนหยุดยาก
ทั้งสามกลไกเป็นเหตุผลที่งานวิจัยนี้เก็บทั้งเวลาที่ใช้ ประเภทเนื้อหา จำนวนวันที่ติดตามข่าวเชิงลบ
และตัวแปรด้านการนอนและการออกกำลังกาย

\subsection{งานวิจัยที่เกี่ยวข้อง}

Orben และ Przybylski \cite{orben2019association} วิเคราะห์ข้อมูลวัยรุ่น 355,358 คน
ด้วยวิธี specification curve analysis และพบว่าการใช้เทคโนโลยีดิจิทัลสัมพันธ์ทางลบกับสุขภาวะ
แต่มีขนาดเล็กมาก คืออธิบายความแปรปรวนของสุขภาวะได้ไม่เกินร้อยละ 0.4
Sala และคณะ \cite{sala2024umbrella} ทบทวนงานทบทวนวรรณกรรม 24 ชิ้น
และสรุปว่าผลของโซเชียลมีเดียขึ้นกับลักษณะผู้ใช้ รูปแบบการใช้งาน และการออกแบบแพลตฟอร์ม
Dominguez-Rodriguez และคณะ \cite{dominguez2025climate}
ศึกษาผู้ตอบ 365 คนในเนเธอร์แลนด์และเยอรมนี
และพบว่าความวิตกกังวลทำนาย doomscrolling ได้อย่างมีนัยสำคัญ
ขณะที่ภาวะซึมเศร้าไม่มีนัยสำคัญในภาพรวม
ส่วน Manzar และคณะ \cite{manzar2025dass21} ตรวจสอบ DASS-21 ในนักศึกษามหาวิทยาลัย 364 คน
พบว่าแบบวัดมีความเชื่อมั่นสูง (แอลฟาของครอนบาค 0.93)
แต่โครงสร้างสามมิติไม่ชัดเจนในทุกกลุ่ม

ผู้วิจัยยังไม่พบงานที่ศึกษาพร้อมกันทั้งชั่วโมงการใช้งานที่อ่านจากค่าที่เครื่องบันทึกไว้
ประเภทเนื้อหาที่เสพ และคะแนนสามมิติของ DASS-21 ในนักศึกษาไทย
งานวิจัยนี้จึงศึกษาประเด็นดังกล่าวในขอบเขตมหาวิทยาลัยเดียว
